%!TEX program = xelatex
\documentclass[cn,hazy,blue,14pt,screen]{elegantnote}
\title{计算机基础知识培训}

\author{小布老师}
\institute{樱桃溪学院}

% \version{2.50}
% \date{\zhdate{2022/12/31}}

\usepackage{array}

\begin{document}

\maketitle

\centerline{
  \includegraphics[width=0.3\textwidth]{ccalogo.png}
}

\newpage
\section{十进制、二进制和十六进制}

\begin{enumerate}[label=\arabic*).]
  \item \textit{何谓十进制？}\\
    逢十进一，数字有十个：0,1,2,3,4,5,6,7,8,9
  \item \textit{何谓二进制？}\\
    逢二进一，数字有两个：0,1
  \item \textit{何谓X进制？}\\
    逢X进一，数字有X个：0,1,...
\end{enumerate}

\newpage
\section{十进制}
\centerline{
  \includegraphics[width=1.0\textwidth]{f0008.pdf}
}

{\Huge
\[
3 \times 100 + 1 \times 10 + 5 = 315
\]
}

\newpage
\section{二进制}
\centerline{
  \includegraphics[width=1.0\textwidth]{f0009.pdf}
}

{\Huge
\[
1 \times 8 + 1 \times 4 + 0 \times 2 + 1 = 13
\]
}

\newpage
\section{十六进制}
\begin{enumerate}
  \item 共计16个数字：0、1、2、3、4、5、6、7、8、9、A、B、C、D、E、F
  \item 不分大小写：a、b、c、d、e、f
\end{enumerate}

\centerline{
  \includegraphics[width=1.0\textwidth]{f0010.pdf}
}

{\Huge
\[
C \times 16 + 5 = 12 \times 16 + 5 = 192 + 5 = 197
\]
}

\section{十进制、二进制和十六进制之间的转换}
\centerline{
  \includegraphics[width=0.8\textwidth]{f0007.pdf}
}

%\printbibliography[heading=bibintoc, title=\ebibname]

\end{document}
